\chapter{Dental Implants}

\section*{Overview}
A dental implant is a biocompatible post placed in the jawbone to support a crown, bridge, or denture after tooth loss. The exact approach, setting, and preparation depend on the indication, anatomy, and the patient's health.

\section*{Why It Is Done}
Implants restore chewing function, support dental prostheses, and help preserve jawbone structure when a patient has suitable bone and oral health.

\section*{How It Is Performed}
After imaging and treatment planning, the implant is placed in the jaw and allowed to integrate with bone. An abutment and final crown or prosthesis are attached after healing.

\section*{Risks and Recovery}
Pain, infection, nerve or sinus injury, failure of bone integration, peri-implant disease, and prosthesis complications are possible. Careful oral hygiene and regular dental follow-up are essential.

\section*{Outcome and Follow-up}
The expected outcome depends on the reason for the procedure, the person's health, and whether additional treatment is needed. The clinician reviews the result with the patient, explains any next steps, and sets follow-up according to the findings and recovery.

\section*{When to Seek Medical Care}
Follow the procedure team's specific instructions. Seek urgent medical assessment for severe or worsening pain, uncontrolled bleeding, fever or signs of infection, trouble breathing, chest pain, fainting, new weakness or confusion, inability to urinate, or any rapidly worsening symptom. The appropriate warning signs depend on the procedure and the patient's medical condition.

\section*{References}
\begin{enumerate}
  \item MedlinePlus. Dental Implants. U.S. National Library of Medicine; 2025.
  \item Current Medical Diagnosis and Treatment. McGraw-Hill; 2025.
\end{enumerate}
\clearpage
